\documentclass[a4paper,12pt]{article}
\usepackage{graphicx, setspace, array, xcolor, mathtools, contour, tikz, amsthm, setspace, amsmath,amsfonts,amssymb, subcaption, fancyhdr, lipsum,multicol, geometry, gensymb}
\usepackage{colortbl}
\geometry{a4paper, margin=0.9in}
\contourlength{0.5pt}
\usetikzlibrary{patterns}
\usepackage[english]{babel}
\renewcommand{\qedsymbol}{$\blacksquare$}
\definecolor{darkgreen}{rgb}{0.0, 0.5, 0.0}
\pagestyle{fancy}
\fancyhf{}
\renewcommand{\headrulewidth}{0.4pt}
\fancyhead[L]{\rightmark}
\fancyhead[R]{\thepage}
\title{Modern Algebra}
\author{Assignment 6\\ \\ Yousef A. Abood\\ \\ ID: 900248250}
\date{November 2025}
\setlength{\parindent}{0pt}
\singlespacing
\parskip=1mm
\setcounter{section}{5}
\setcounter{subsection}{1}
\onehalfspacing
\begin{document}
\maketitle
\noindent\makebox[\linewidth]{\rule{15cm}{0.4pt}}
\section{Group Isomorphisms}
% 1, 2, 4, 8, 11, 12, 14, 21, 22, 24, 34.
\subsection*{Problem 1}
Observe that $\phi : Z \to 2Z$ such that $ \phi (x)=2x$ is a bijection: \\
Suppose $\phi(a)=\phi(b) \implies 2a=2b \implies a=b.$ Thus, the function is clearly injective.\\
The function is clearly surjective, as for every $t$ in the codomain, by definiton, $t=2k$ where $k \in \mathbb{Z}.$ Clearly, $\phi(k)=t.$
It is also a homomorphism as
\[\phi(a+b)=2(a+b)=2a+2b=\phi(a)+\phi(b)\]
\subsection*{Problem 2}
By Theorem 6.2.2. in the Lecture Notes, if there exist an isomorphism between two cyclic groups, then this isomorphism has to send every generator in the domain to a generator in the codomain. We know that $\mathbb{Z}$ with addition is cyclic and it has only two generators $1, -1$. So we have only two possible isomorphisms from $\mathbb{Z}$ to itself. These two automorphisms are $\phi(x)=x, g(x)=-x$. They are clearly injective and serjective.
Observe:
\begin{align*}
    &\phi(a+b)=a+b=\phi(a)+\phi(b)\\
    &g(a+b)=-(a+b)=-a-b=g(a)+g(b).
\end{align*}
Both functions are homomorphisms. Hence, $\phi, g$ are the only possible automorphisms and $Aut(\mathbb{Z})={\phi, g}.$ 
\subsection*{Problem 4}
By Theorem 6.2.5. in Lecture Notes, if we have an isomorphism between two groups $f: G \to H$, where one of these of groups is cyclic, then the other group must be cyclic. We observe that $\langle 3 \rangle=U_{10}$ but $U_{8}$ is not cyclic. Thus, no isomorphism can be found between $U_{10}, U_8.$
\subsection*{Problem 8}
\begin{proof}
    We will call the mapping $a\to \log_{10}{a}$ by $g$. \\
    First, we need to show that this function is bijective. We start by showing injectivity. Suppose $g(x)=g(y)$, then $\log_{10}x=\log_{10}y$. We know that the logarithmic function has an inverse, so $x=y$ and $g$ is injective. Next, we show serjectivity. By definition of the codomain, we can see that every element in the codomain can be written as $\log_{10}t=k$, where $t,k\in \mathbb{R^+}, \mathbb{R}$ respectively. Observe that $10^k=t$ and $g(t)=log_{10}{10^k}=k.$ Thus, the function is serjective. Therefore, we showed that $g$ is a bijection from $\mathbb{R^+}\to \mathbb{R}.$
\\  Second, we will show that $g$ is a homomorphism. Observe
\begin{align*}
    g(xy)=\log_{10}{xy}=\log_{10}{x}+\log_{10}{y} \text{  (Using properties of the logarithmic function).}
\end{align*}
Clearly a homomorphism. Therefore, we showed that the mapping $a\to \log_{10}{a}$ is an isomorphism from $\mathbb{R^+}$ under multiplication to $\mathbb{R}$ under addition.
\end{proof}
\subsection*{Problem 11}
\begin{align*}
    \phi(a^3b^{-2})&=\phi(a^3)+\phi(b^{-2})\\
    &=\phi(a \cdot a \cdot a)+\phi(b^{-1}\cdot b^{-1})\\
    &=\phi(a)+\phi(a)+\phi(a)+\phi(b^{-1})+\phi(b^{-1})\\
    &=3\bar{a}+\phi(b^{-1})+\phi(b^{-1})=3\bar{a}+2\phi(b^{-1}).
\end{align*}
Since the isomorphism maps inverses to inverses, then $\phi(b^{-1})=-\phi(b).$
Then the final expression is $3\bar{a}-2\bar{b}.$ 
\subsection*{Problem 12}
\begin{proof}
    Assume $\alpha$ is an automorphism, we need to show that $G$ is abelien. Observe since $\alpha$ is a homomorphism, then pick $s,t \in G$:
    \begin{align*}
        \alpha(st)=(st)^{-1}=t^{-1}s^{-1}\\
        \alpha(s)\alpha(t)=s^{-1}t^{-1}.
    \end{align*}
    And we know that $\alpha(st)=\alpha(s)\alpha(t)$, then $t^{-1}s^{-1}=s^{-1}t^{-1}.$ Since, $s,t \in G$ then $s^{-1},t^{-1} \in G.$ then the group is abelien.
    For the backward direction, assume the group is abelien, we need to show that $\alpha$ is an automorphism. First, we show it is bijective. 
    Since, $G$ is a group, then every element has a unique inverse and $\alpha$ is clearly both injective and serjective. That is, if $\alpha(a)=\alpha(b)$, then $a=b$ as inverses are unique, so $\alpha$ is injective. And since for an element in the codomain $a^{-1}$, it is also in the group and has a unique inverse $a$ and $\alpha$ is surjective. Second, we show $\alpha$ is a homomorphism, pick $s,t \in G$ and observe:
    \begin{align*}
        \alpha(st)=(st)^{-1}=t^{-1}s^{-1}=s^{-1}t^{-1}=\alpha(s)\alpha(t).
    \end{align*}
    Hence, $\alpha$ is a homomorphism.\\
    Therefore, $\alpha(g)=g^{-1}$ for all $g \in G$ is an isomorphism iff $G$ is abelien.
\end{proof}
\subsection*{Problem 14}
Observe $\mathbb{Z}_1, \mathbb{Z}_2$ are not isomorphic, as $|\mathbb{Z}_1| \ne |\mathbb{Z}_2|.$
By theorem 6.3.5. in Lecture Notes, we know that $Aut(\mathbb{Z}_n)\cong U_n.$ So, $Aut(\mathbb{Z}_2) \cong U_2 = {1}$ and $Aut(\mathbb{Z}_1) \cong U_1= {1}.$ Since, $U_1=U_2$ then $Aut(\mathbb{Z}_2) \cong U_1$. By transativity, $Aut(\mathbb{Z}_1) \cong Aut(\mathbb{Z}_2).$
\subsection*{Problem 21}
\begin{proof}We see that $H$ contains the permutations that fix $1$ and permute the other elements. Since, $\beta \in S_5$ then the remaining elements are $\{2,3,4,5\}$, which are 4 elements. We can see that is isomorphic to $S_4$. Similarly, the group $K$ has the permutations that fix $2$ and permute $\{1,3,4,5\}$ which is also isomorphic $S_4$. Thus, $H \cong S_4 \wedge K \cong S_4 \implies H \cong K.$
We use the same approach if $\beta \in S_n$ where $n\ge 3$. We know $\beta$ fixes an element and permute the others, then the group is isomorphic to $S_{n-1}.$
\end{proof}
\subsection*{Problem 22}
\textit{This is just a generalization to what we did in Problem 1.}
\begin{proof}
    Observe that $\phi : \mathbb{Z} \to k\mathbb{Z}$ such that $ \phi (x)=kx$ is a bijection, where $k \in \mathbb{Z^{*}}$: \\
Suppose $\phi(a)=\phi(b) \implies ka=kb \implies a=b$, where $a,b \in \mathbb{Z}.$ Thus, the function is clearly injective.\\
The function is clearly surjective, as for every $t$ in the codomain, by definiton, $t=km$ where $m \in \mathbb{Z}.$ Clearly, $\phi(m)=t.$
It is also a homomorphism as
\[\phi(a+b)=k(a+b)=ka+kb=\phi(a)+\phi(b)\]
Hence, we know that $K\mathbb{Z}$ is a subgroup of $\mathbb{Z}$ for all $k \in \mathbb{Z}^*$ and we showed it is isomorphic to $\mathbb{Z}.$
\end{proof}
\subsection*{Problem 24}
\begin{proof}
    Since $H$ is the set of the elements fixed by the automorphism $\phi$. We can see that the identity element is in $H$ as the automorphism must send the identity to itself. Next, we show it is closed under the operation. Pick $a,b \in H$, we want to show that $ab \in H.$ Since $\phi$ is a homomorphism, then $\phi(ab)=\phi(a)\phi(b).$ But we know that since $a,b \in H$, then $\phi$ fixes them. Thus, $\phi(ab)=\phi(a)\phi(b)=ab$ and $ab \in H.$ \\
    Last, we proof that $H$ is closed under taking inverses. Since $\phi$ is an isomorphism, then $\phi(a^{-1})=\phi(a)^{-1}=(\phi(a))^{-1}=a^{-1}.$ Hence, it is closed under taking inverses.
    \\Therefore, we showed that $H$ is a subgroup of $G.$
\end{proof}
\subsection*{Problem 34}
\begin{proof}
    Since $K$ is a subgroup of $G$, then it has the identity element. We know that $\phi$ is an isomorphism from $G \to \bar{G}$ and it must send the identity element in $G$ to the identity in $\bar{G}$. Thus, we showed $\phi(K)$ is not empty and it has the identity element.
    Pick elements $a,b \in \phi(K)$ such that $a=\phi(s), b=\phi(t)$, where $s,t \in K.$ Observe that $ab=\phi(s)\phi(t)=\phi(st)$ and $st \in K$ since $K$ is a subgroup. Hence, $\phi(K)$ is closed under the operation.\\
    Next, we show that $\phi(K)$ is closed under taking inverses. See that $\phi(s^{-1})=\phi(s)^{-1}=a^{-1}$, so every element in $\phi(K)$ has an inverse and $\phi(K)$ is closed under taking inverses.
    Therefore, we proved that $\phi(K)$ is a subgroup of $\bar{G}$ if $K$ is a subgroup of $G$, where $\phi$ is an isomorphism from $G$ to $\bar{G}.$
\end{proof}

\end{document}